\chapter{Challenges and Limitations}

\section{Mean Flow Training Instability}

The most significant challenge encountered was MF training divergence. The loss
rose from epoch 1 and never recovered. Several factors likely contributed:

\begin{itemize}
  \item \textbf{Learning rate.} The submitted configuration uses
        $5\!\times\!10^{-5}$, four times lower than FM and four times lower than
        the default used in the Mean Flow paper~\cite{geng2025mean}. This slows
        convergence and may interact poorly with the JVP-based target.
  \item \textbf{Small batch size.} Batch 32 introduces noisy gradient estimates
        for the JVP target, which itself depends on the current network output.
        Larger batches could stabilise the target.
  \item \textbf{Short training.} 30 epochs at batch 32 processes approximately
        1.5 million samples, versus 4.99 million for FM at batch 128 / 100 epochs.
        The network may need significantly more gradient steps to learn a stable
        mean-velocity field.
  \item \textbf{Additive $r$-conditioning.} The $r$-embedding adds a 3-channel
        offset to the input tensor. A stronger conditioning mechanism (e.g.,
        concatenation or cross-attention) may be needed for the backbone to
        disentangle the two time endpoints.
  \item \textbf{No EMA or gradient clipping.} The paper uses EMA and gradient
        clipping to stabilise JVP-based training; neither is present in the
        submitted configuration.
\end{itemize}

\section{JVP Numerical Sensitivity}

The JVP computation is sensitive to the magnitude of the tangent vector
$v_\theta(z_t, t)$. Early in training, when $v_\theta$ is uninitialised noise,
the JVP output can be very large, producing unstable targets. Without gradient
clipping, the first few epochs may inject large loss spikes that compound through
the rest of training.

\section{Unmatched Compute Budgets}

The three submitted runs differ in batch size, epoch count, learning rate, and
GPU hardware context (FM-LN was trained in a different session on the same
machine). Attributing quality differences to the algorithm alone is therefore
not valid. The MF vs FM comparison in particular conflates training budget
differences with objective differences.

\section{Evaluation Limitations}

\begin{itemize}
  \item FID and IS are estimated from 1{,}000 generated samples. Standard
        benchmarks use 50{,}000; estimates from 1{,}000 samples have substantially
        higher variance and cannot be compared to published results.
  \item IS does not compare against real images and does not detect memorisation
        or mode collapse; it rewards only confidence and diversity of Inception
        classifier predictions.
  \item The FID reference is computed from 1{,}000 real test images rather than
        the full test set, introducing additional variance.
\end{itemize}

\section{Timing and Hardware Attribution}

Logged sampling times differ across runs because measurements were taken in
different GPU sessions with different system loads. The absence of a complete
hardware fingerprint (GPU model, driver, TGP, clock state, background processes)
means that absolute cross-run latency differences cannot be attributed to the
algorithm.

\section{Missing Early Checkpoints}

Checkpoints are saved every ten epochs; no trained images exist for epochs 1--9
in any run. The training history UI therefore begins at epoch 10. This gap
prevents studying early-training behaviour qualitatively.

\section{Documentation Drift}

Earlier root README text described placeholders after working implementations
were present. This drift risked misrepresenting project status to external
readers. The publication workflow therefore added updated setup and training
instructions, scoped ignore rules that retain source packages, and clean-clone
Windows/Linux validation. The incident remains documented here as an engineering
lesson: implementation and public reproducibility guidance must be reviewed as
one deliverable.
